\documentclass[10pt,a4paper]{article}

\usepackage[utf8]{inputenc}
\usepackage[T1]{fontenc}
\usepackage[brazilian]{babel}
\usepackage[a4paper,top=0.8cm,bottom=0.8cm,left=1.4cm,right=1.4cm]{geometry}
\usepackage[hidelinks]{hyperref}
\usepackage{titlesec}
\usepackage{enumitem}
\usepackage{xcolor}
\usepackage{parskip}
\usepackage{array}
\usepackage{fontawesome5}

\definecolor{accent}{HTML}{2E3440}
\definecolor{muted}{HTML}{6B7280}

\pagestyle{empty}
\setlength{\parindent}{0pt}
\setlength{\parskip}{1pt}

\titleformat{\section}
  {\color{accent}\normalsize\bfseries}
  {}{0pt}{\MakeUppercase}[\vspace{-1pt}\color{accent}\titlerule]
\titlespacing*{\section}{0pt}{5pt}{4pt}

% In-column section header (for use inside minipages where \section* layout breaks)
\newcommand{\colsection}[1]{%
  \par\vspace{1pt}%
  {\color{accent}\normalsize\bfseries\MakeUppercase{#1}}\par\nobreak
  \vspace{-3pt}\noindent\textcolor{accent}{\rule{\linewidth}{0.4pt}}\par\vspace{2pt}%
}

\setlist[itemize]{leftmargin=1.1em,itemsep=0pt,topsep=1pt,parsep=0pt}

\newcommand{\role}[4]{%
  \par\vspace{1pt}%
  \noindent\begin{tabular*}{\linewidth}[t]{@{}l@{\extracolsep{\fill}}r@{}}
    \textbf{#1} \textit{\textcolor{muted}{#2}} & \textcolor{muted}{#3} \\
  \end{tabular*}\par\nobreak
  \vspace{1pt}#4\par
}

\newcommand{\edu}[3]{%
  \par\vspace{1pt}%
  \noindent\begin{tabular*}{\linewidth}[t]{@{}l@{\extracolsep{\fill}}r@{}}
    \textbf{#1} \textit{\textcolor{muted}{#2}} & \textcolor{muted}{#3} \\
  \end{tabular*}\par
}

\newcommand{\project}[4]{%
  \par\vspace{1pt}%
  \noindent\begin{tabular*}{\linewidth}[t]{@{}l@{\extracolsep{\fill}}r@{}}
    \textbf{#1} \textcolor{muted}{--- #2} & \href{https://#3}{\textcolor{muted}{#3}} \\
  \end{tabular*}\par\nobreak
  #4\par
}

\newcommand{\contact}[2]{\textcolor{muted}{#1}~\,#2}

\begin{document}

%----- Cabeçalho: nome/cargo (esquerda) | contato (direita) -----
\noindent
\begin{minipage}[b]{0.55\linewidth}
  {\Huge\bfseries André Brandão}\\[3pt]
  {\color{muted}\normalsize Engenheiro de Plataforma \& Backend}
\end{minipage}%
\hfill
\begin{minipage}[b]{0.42\linewidth}
  \small
  \setlength{\tabcolsep}{4pt}
  \renewcommand{\arraystretch}{1.15}
  \begin{tabular}{@{}l@{\hspace{12pt}}l@{}}
    \contact{\faEnvelope}{\href{mailto:eu@andrebrandao.dev}{eu@andrebrandao.dev}} &
      \contact{\faGithub}{\href{https://github.com/andre-brandao}{github.com/andre-brandao}} \\
    \contact{\faGlobe}{\href{https://andrebrandao.dev}{andrebrandao.dev}} &
      \contact{\faMapMarker*}{Belo Horizonte, Brasil} \\
    \multicolumn{2}{@{}l@{}}{\contact{\faLanguage}{Português (nativo) \textbullet{} Inglês (proficiência profissional)}} \\
  \end{tabular}
\end{minipage}

\vspace{3pt}

%----- Perfil | Habilidades (lado a lado) -----
\noindent
\begin{minipage}[t]{0.46\linewidth}
  \colsection{Perfil}
  Engenheiro de plataforma e backend que entrega ponta a ponta, da AWS e Terraform
  no trabalho a um homelab declarativo em NixOS + Proxmox em casa. Desenvolvedor
  Full Stack na A3Data, liderando o A3Copilots, plataforma interna de IA que
  expõe os sistemas da empresa a agentes via servidores MCP. Cofundei a Crossgeo,
  onde construí o sistema de gestão de uma distribuidora de bebidas e uma
  plataforma de entregas de última milha, ambos em produção. Bacharel em
  Ciência da Computação pela PUC-MG (2026).
\end{minipage}%
\hfill
\begin{minipage}[t]{0.52\linewidth}
  \colsection{Habilidades}
  \small
  \begin{tabular}{@{}p{0.25\linewidth}p{0.71\linewidth}@{}}
    \textbf{Linguagens} & TypeScript, Python, Go, C, Java, Rust, Bash \\
    \textbf{Backend} & APIs REST / gRPC, microsserviços, orientado a eventos
      (SQS, EventBridge, DynamoDB Streams), ECS, Lambda, RDS, DynamoDB, S3, CloudFront \\
    \textbf{Infra} & AWS, Cloudflare, Terraform, Pulumi/SST, Backstage,
      Docker, Proxmox, GitHub Actions, Vault, IAM, Linux, Nix / NixOS, Git, SOPS \\
    \textbf{Dados} & PostgreSQL, SQLite, Redis, SQL, NoSQL, Airflow, Spark \\
    \textbf{Observabilidade} & OpenTelemetry, Prometheus / Grafana, tracing distribuído \\
    \textbf{IA \& LLMs} & Amazon (Bedrock/AgentCore), sistemas multiagentes, servidores MCP, Langchain \\
  \end{tabular}
\end{minipage}

\section*{Experiência}
\role{Desenvolvedor Full Stack}{A3Data --- Belo Horizonte, Brasil}{06/2026 -- Atual}{
  \begin{itemize}
    \item Desenvolvedor principal do \textbf{A3Copilots}, plataforma interna de IA
          que conecta os sistemas da empresa (CRM Pipedrive, Google Drive, AWS
          Partner Network, catálogo do Backstage) a agentes via servidores MCP
          com controle de acesso por papel (RBAC), de forma que cada agente só
          vê os dados que o usuário tem permissão de acessar.
    \item Projetei a estrutura do projeto e construí do zero o CI/CD e a
          infraestrutura, incluindo ambientes de preview efêmeros por pull
          request, para que toda mudança seja revisada em um deploy real antes
          do merge.
    \item \textit{Stack:} TypeScript, Node, MCP, Backstage, Terraform, AWS, GitHub Actions.
  \end{itemize}
}
\role{Estagiário em Engenharia de ML}{A3Data --- Belo Horizonte, Brasil}{03/2025 -- 06/2026}{
  \begin{itemize}
    \item Implementei a camada de observabilidade do sistema multiagente de
          cotações do Grupo Elfa no Amazon Bedrock AgentCore, a primeira
          implantação do AgentCore em produção na América Latina: tracing
          distribuído, KPIs de negócio e dashboards via OpenTelemetry para mais
          de 800 mil cotações por mês.
          \href{https://aws.amazon.com/pt/blogs/aws-brasil/grupo-elfa-aplicando-amazon-bedrock-agentcore-para-observabilidade-de-sistema-multiagente-em-escala/}{Publicado no blog de engenharia da AWS Brasil}.
    \item Construí a plataforma interna de desenvolvedores em Backstage com um
          motor de templates para agentes de WhatsApp e LLM web, pipelines de
          ETL (Airflow + Spark) e projetos de ML, cada um já com infra em
          Terraform, pipelines de deploy, monitoramento e IAM.
    \item Provisionei e implantei agentes e serviços em produção na AWS em ECS,
          Lambda, DynamoDB (incl.\ Streams) e RDS inteiramente via Terraform.
    \item \textit{Stack:} Python, Terraform, Backstage, OpenTelemetry, Airflow, Spark, AWS.
  \end{itemize}
}
\role{Desenvolvedor Full Stack, Cofundador}{Crossgeo --- Belo Horizonte, Brasil}{06/2024 -- 09/2026}{
  \begin{itemize}
    \item \textbf{Duarte Distribuidora:} sistema de gestão em produção em uma
          distribuidora de bebidas, cobrindo estoque, vendas, clientes e
          fornecedores. Svelte, TypeScript, SQLite.
    \item \textbf{Rota 88 Delivery:} plataforma de logística de última milha em
          produção. Ingere pedidos do iFood e do OpenDelivery via webhooks em Lambda,
          roda a aplicação web do operador em ECS com WebSockets para despacho
          em tempo real, e entrega um app React Native para os motoboys com
          heurísticas configuráveis de despacho automático.
          Svelte, TypeScript, React Native, PostgreSQL no RDS, AWS.
  \end{itemize}
}
\role{Desenvolvedor Full Stack}{Crossvirus (startup PUC-MG) --- Belo Horizonte, Brasil}{03/2024 -- 2025}{
  Desenvolvi uma plataforma de monitoramento epidemiológico (Dengue, COVID-19)
  usada por duas cidades de MG e coordenei dois estagiários. Astro, Svelte, PostgreSQL.
}

\section*{Projetos}
\project{Homelab}{Nix / NixOS, Terraform, Proxmox, SOPS}{github.com/andre-brandao/nixos}{
  Infraestrutura totalmente declarativa para minhas estações de trabalho (NixOS +
  home-manager) e um homelab em Proxmox rodando Gitea, Vault, Prometheus, Grafana
  e um gateway de LLM. Cada host é um módulo NixOS com um stack Terraform
  correspondente; segredos criptografados com SOPS.
}

\project{GPUi Shell}{Rust, Wayland, D-Bus, IPC, Nix}{github.com/andre-brandao/gpui-shell}{
  Shell de desktop para Wayland acelerado por GPU, em Rust sobre o framework GPUI
  do Zed. Substitui waybar, rofi e swaync por um único app para Hyprland e Niri.
}

\section*{Certificações}
AWS Partner --- Agentic AI Essentials \textbullet{} Generative AI (Technical) \textbullet{}
Digital Sovereignty on AWS (Technical) \textbullet{}
AWS Academy Graduate --- Cloud Foundations.

\section*{Formação}
\edu{Bacharelado em Ciência da Computação}{PUC-MG, Belo Horizonte, Brasil}{2022 -- 2026}

\end{document}
